\chapter[Classical Smooth Functionals]{Classical Smooth Information Functionals}\label{ch:sec:classical-smooth-information-functionals}

This chapter applies the universal bounds to the classical smooth information functionals---Shannon entropy, R\'enyi entropy, $g$-entropy---and derives explicit Lipschitz constants for each. These constants will be compared with the piecewise-smooth bounds of Chapter~\ref{ch:sec:classical-smooth-information-functionals} and used in the applications of Chapter~\ref{ch:sec:applications-of-the-wis-framework}.


\section{Introduction}\label{ch:sec:introduction-smooth_functionals}
The previous chapters established geometric results that apply to
arbitrary smooth scalar functionals and, more specifically, to Bregman
divergences. This chapter demonstrates how several classical
information-theoretic quantities naturally fit into the WIS framework.
The emphasis is not on re-deriving well-known formulas, but on showing
that these quantities become instances of the same geometric theory once
expressed on the WIS manifold.  For background on classical information
measures, see \cite{cichocki2010families, sason2016fdivergence, reid2011information}.
\subsection{Shannon Entropy}\label{ch:cor:16-1}
The Shannon entropy \cite{cover2012elements} is
\(H(P)=-\sum_i p_i\log p_i.\)
Since \(H\) is smooth on the interior of the probability simplex, it
defines a smooth scalar field on the WIS manifold.
\begin{corollary}[Universal Stability]

If the WIS gradient of \(H\) is bounded on a region
\(\Omega\subseteq\mathcal P\), then
$\|H(P)-H(Q)\| \le L_H\sqrt{\mathcal D(P,Q)}.$
Near any stationary point of \(H\),
\[ \|H(P)-H(P^{\star})\| \le
C_H\mathcal D(P,P^{\star}). \]
These bounds follow directly from Chapter~\ref{ch:sec:information-functionals-on-the-wis-manifold}.
\end{corollary}
\section{Kullback--Leibler Divergence}\label{ch:sec:kullback--leibler-divergence}
For a fixed reference distribution \(Q\),
\[ D_{\mathrm{KL}}(P\|Q) = \sum_i p_i
\log\frac{p_i}{q_i}. \]
Viewed as a function of \(P\), the KL divergence is generated by the
convex potential
$\phi(x)=\sum_i e^{x_i},$
up to the coordinate normalization induced by the WIS embedding.
Consequently, the results of Chapter~\ref{ch:sec:universal-geometric-bounds} apply immediately.
\begin{corollary}[KL Average Hessian]
\label{ch:cor:16-2}
There exists an Average Hessian Operator
$\overline H_{\mathrm{KL}}(P,Q)$
such that
\[ D_{\mathrm{KL}}(P\|Q) = \frac12 \Delta^T
\overline H_{\mathrm{KL}} \Delta. \]
Furthermore,
\[ \frac{\underline\lambda}{2} \mathcal D \le
D_{\mathrm{KL}} \le \frac{\overline\lambda}{2}
\mathcal D, \]
whenever the hypotheses of Theorem~\ref{ch:thm:15-3} are satisfied.
\end{corollary}
\subsection{R\'enyi Entropy}\label{ch:sec:renyi-entropy}
For \(\alpha>0\), \(\alpha\neq1\),
\[ H_\alpha(P) = \frac1{1-\alpha} \log
\sum_i p_i^{\alpha}. \]
Since the interior of the simplex excludes vanishing probabilities,
The R\'enyi entropy is smooth on the interior simplex.
Therefore the universal bounds of Chapter~\ref{ch:sec:information-functionals-on-the-wis-manifold} apply without modification,
yielding
\[
|H_\alpha(P)-H_\alpha(Q)|
\le
L_{\alpha,K}\,\sqrt{\mathcal D(P,Q)}
\]
on any compact subset \(K\subset\Delta_n^\circ\).
We now derive an explicit expression for the Lipschitz constant
\(L_\alpha\) and compare the resulting estimate with known bounds from
the literature.
\section{Gradient of the R\'enyi Entropy}\label{ch:sec:gradient-of-the-renyi-entropy}
Recall (Chapter~\ref{ch:sec:pullback-riemannian-metric}) that the pullback metric on the posterior manifold is induced by
the embedding \(F=Q\circ L\). In logarithmic coordinates, the Riemannian
gradient of a smooth functional \(S\) satisfies
\[
g(\nabla_g S, v) = dS(v) \qquad\text{for all }v.
\]
For the R\'enyi entropy of order \(\alpha>0\), \(\alpha\neq1\),
$H_\alpha(P)=\frac1{1-\alpha}\log\sum_i P_i^{\alpha},$
the differential is
\[
dH_\alpha(P)(v)
=
\frac{\alpha}{1-\alpha}
\,
\frac{\sum_i P_i^{\alpha-1}v_i}{\sum_i P_i^{\alpha}}.
\]
In the WIS logarithmic coordinates \(a_i=\log(nP_i)\), the squared norm
of the gradient is
\[
\|\nabla_g H_\alpha(P)\|_g^2
=
\frac{\alpha^2}{(1-\alpha)^2}
\,
\sum_{i=1}^n
\Bigl(
\frac{P_i^{\alpha}}{\sum_j P_j^{\alpha}}
-
P_i
\Bigr)^2,
\]
because the projection of the logarithmic gradient onto the tangent
space of the simplex removes the shift direction and the resulting
components sum to zero.  A convenient explicit upper bound, used for
the constants below, is obtained by bounding each summand:
\[
\|\nabla_g H_\alpha(P)\|_g^2
\le
\frac{\alpha^2}{(1-\alpha)^2}
\,
\frac{
\sum_i P_i^{2\alpha-2}
}{
\bigl(\sum_i P_i^{\alpha}\bigr)^2
}
\cdot
\max_i P_i.
\]
Both expressions are continuous on the interior and bounded on any
compact subset \(K\). By Theorem~\ref{ch:thm:14-1}, the Lipschitz estimate
follows with \(L_{\alpha,K}\) equal to the supremum of the exact
gradient norm over \(K\); the upper bound above yields an explicit
constant that is valid (though not tight) on \(K\).
\section{Improved Bound via the Stationary Point}\label{ch:sec:improved-bound-via-the-stationary-point}
The R\'enyi entropy for \(\alpha>0\) has a unique stationary point at the
uniform distribution \(U_n\). Expanding around this stationary point and
applying Theorem~\ref{ch:thm:14-2} yields a \textbf{quadratic} bound
\[
|H_\alpha(P)-H_\alpha(U_n)|
\le
C_{\alpha,K}\,\mathcal D(P,U_n),
\]
valid on any compact subset \(K\) containing \(U_n\). The constant
\(C_{\alpha,K}\) is half the operator norm of the WIS Hessian of
\(H_\alpha\) on \(K\).
For \(\alpha>1\), the Hessian norm is maximized at the boundary of the
simplex, giving an explicit bound
\[
C_{\alpha,K}^2
\le
\frac{\alpha^2}{(\alpha-1)^2}
\cdot
\frac{
\sum_i P_i^{2\alpha-2}
}{
\bigl(\sum_i P_i^{\alpha}\bigr)^2
}
\cdot
\frac{\max_i P_i}{\min_i P_i},
\]
which on a neighbourhood where \(\min_i P_i\ge\varepsilon>0\) and
\(n\le N\) simplifies to
\[
C_{\alpha,K}
\le
\frac{\alpha}{|\alpha-1|}
\cdot
\frac{1}{\varepsilon^{\alpha}\,N^{\alpha/2}}.
\]
\section{Comparison with Known Bounds}\label{ch:sec:comparison-with-known-bounds}
Let \(U_n\) denote the uniform distribution on \(n\) symbols. The
R\'enyi divergence of order \(\alpha\) is
\[
D_\alpha(P\|U_n)=\frac1{\alpha-1}\log\sum_i n^{\alpha-1}P_i^{\alpha}.
\]
A standard identity \cite{vanerven2014renyi} relates the R\'enyi
entropy to the R\'enyi divergence from the uniform distribution:
\[ H_\alpha(P) = \log_2 n - D_\alpha(P\|U_n). \tag{16.1}\label{eq:renyi-identity} \]
This identity is exact. It provides a complete description of the
R\'enyi entropy once the R\'enyi divergence is known, and is therefore
strictly tighter than any bound expressed in terms of a different
metric. The present geometric bounds do not compete with~\eqref{eq:renyi-identity} when the
R\'enyi divergence is explicitly computable.
The contribution of the geometric approach is of a different nature.
The UOD \(\mathcal D(P,U_n)\) is a purely geometric quantity---the
squared Euclidean distance between the logarithmic embeddings of \(P\)
and \(U_n\) in the quotient space. It can be computed directly from the
posterior probabilities without evaluating any parametric family of
divergences. The geometric bounds
\[
|H_\alpha(P)-\log_2 n| \le L_\alpha\sqrt{\mathcal D(P,U_n)},
\qquad
|H_\alpha(P)-\log_2 n| \le C_\alpha\,\mathcal D(P,U_n)
\]
therefore provide an \textbf{independent} estimate of the R\'enyi entropy
solely from the geometry of the posterior manifold. They are most useful
in settings where the R\'enyi divergence is difficult to compute but the
UOD is accessible, for instance when the logarithmic embedding is the
natural coordinate system of the problem.
Moreover, the Lipschitz and quadratic bounds apply to \textbf{any} smooth
functional on the posterior manifold, including Shannon entropy,
collision entropy, and security measures such as min-entropy and Bayes
vulnerability. The R\'enyi entropy is one instance of a general
principle rather than a specialized estimate.
\begin{remark}[Tightness of the Quadratic Bound for Small Defects]
For small defects \(\mathcal D(P,U_n)\ll 1\), the quadratic bound
$|H_\alpha(P)-\log_2 n| \le C_{\alpha,K}\,\mathcal D(P,U_n)$
is nearly tight: the constant \(C_{\alpha,K}\) coincides with half the
spectral norm of the WIS Hessian, which for the uniform distribution
attains the exact value
\[
C_{\alpha,U_n}
=
\frac{\alpha}{2|\alpha-1|}\cdot\frac{1}{n}.
\]
In this regime the geometric bound provides a numerically accurate
estimate, whereas the exact identity~\eqref{eq:renyi-identity} requires computing the
full R\'enyi divergence.
\end{remark}

\begin{remark}[Improved Bound via Exact Hessian]
\label{rem:improved-hessian}

The quadratic bound can be sharpened by computing the WIS Hessian of the
R\'enyi entropy explicitly at every interior point, rather than
bounding its norm on a compact set.  The Hessian in logarithmic
coordinates is

\[
\operatorname{Hess}_g H_\alpha(P)
=
\frac{\alpha}{(1-\alpha)}
\,
\frac{
\operatorname{diag}(P^{\alpha-2})
}{
\sum_i P_i^{\alpha}
}
-
\frac{\alpha^2}{(1-\alpha)^2}
\,
\frac{
P^{\alpha-1}(P^{\alpha-1})^T
}{
\bigl(\sum_i P_i^{\alpha}\bigr)^2
},
\]

where the first term is a diagonal matrix and the second is a rank-one
correction.  The operator norm of this Hessian satisfies

\[
\|\operatorname{Hess}_g H_\alpha(P)\|
=
\frac{\alpha}{|\alpha-1|}
\,
\frac{
\max_i P_i^{\alpha-2}
}{
\sum_i P_i^{\alpha}
},
\]

which is attained in the direction of the coordinate with the smallest
probability.  Consequently, the optimal quadratic constant at a point
\(P\) is

\[
C_\alpha^{\star}(P)
=
\frac{\alpha}{2|\alpha-1|}
\,
\frac{
\max_i P_i^{\alpha-2}
}{
\sum_i P_i^{\alpha}
}.
\]

For the uniform distribution \(U_n\), this reduces to the previously
stated value \(C_{\alpha,U_n}=\alpha/(2|\alpha-1|)\cdot 1/n\).  For a
non-uniform \(P\), the bound~\eqref{eq:renyi-identity} is still tighter when the
R\'enyi divergence is computable, but the pointwise Hessian bound
\(C_\alpha^{\star}(P)\) is sharper than the uniform compact-set bound
\(C_{\alpha,K}\) when \(P\) is far from the uniform distribution.
\end{remark}

\begin{remark}[Interpolated Bound]
\label{rem:interpolated-bound}

The Lipschitz and quadratic bounds can be combined into a single
estimate that is never worse than either individually.  For any
\(P\in\Delta_n^\circ\),

\[
|H_\alpha(P)-\log_2 n|
\le
\min\left\{
L_{\alpha,K}\sqrt{\mathcal D(P,U_n)},
\;
C_{\alpha,K}\,\mathcal D(P,U_n)
\right\}.
\]

The quadratic term dominates for \(\mathcal D(P,U_n)\ll 1\), while the
Lipschitz term dominates for larger defects.  The crossover point is
\(\mathcal D(P,U_n) = (L_{\alpha,K}/C_{\alpha,K})^2\), below which the
quadratic estimate is tighter.  Using the exact Hessian constant
\(C_\alpha^{\star}(P)\) in place of \(C_{\alpha,K}\) shifts this
crossover further toward small defects, ensuring that the quadratic
bound is always at least as tight as the Lipschitz bound in a
neighbourhood of the uniform distribution.
\end{remark}
\begin{example}[Binary Min-Entropy in Bits]
\label{ch:ex:16-6}
Consider a binary distribution \(P=(p,1-p)\) on \(n=2\) symbols, with
uniform reference \(U_2=(1/2,1/2)\). The UOD is
\[
\mathcal D(P,U_2)
=
\bigl(\log(2p)-\bar a\bigr)^2
+
\bigl(\log(2(1-p))-\bar a\bigr)^2
=
\frac12\Bigl(\log(2p)-\log(2(1-p))\Bigr)^2,
\qquad
\bar a=\frac{\log(2p)+\log(2(1-p))}{2},
\]
where all logarithms are taken in base 2. The min-entropy in bits is
$H_\infty(P)=-\log_2\max\{p,1-p\},$
and the uniform reference gives \(H_\infty(U_2)=\log_2 2 = 1\) bit.
For \(p\in[0.1,0.9]\), the Lipschitz constant \(L_\infty\) attains its
maximum at \(p=0.1\) (or symmetrically \(p=0.9\)), giving
\[
L_\infty^2
=
\frac{\max\{p,1-p\}}{(\min\{p,1-p\})^2}
=
\frac{0.9}{(0.1)^2}=90,
\qquad
L_\infty\approx9.49.
\]
Therefore, for any binary distribution with \(p\in[0.1,0.9]\),
\[
|H_\infty(P)-1|
\le
9.49\,\sqrt{\mathcal D(P,U_2)}\quad\text{(bits)}.
\]
In particular, if \(\mathcal D(P,U_2)\le0.01\) (a small defect), then
\[
|H_\infty(P)-1|\le0.949\text{ bits},
\]
meaning the min-entropy deviates from its uniform value of 1 bit by at most
\(0.949\) bits.
If instead \(\mathcal D(P,U_2)\le1\), the bound becomes
\[
|H_\infty(P)-1|\le9.49\text{ bits},
\]
which is vacuous because the min-entropy of a binary distribution is
always between \(0\) and \(1\) bit. The Lipschitz bound is therefore
most informative for small defects, where the linear approximation
is accurate and the resulting estimate is tighter than the trivial
range \([0,1]\).
Concretely, for the bound to be non-vacuous we require
\[
9.49\,\sqrt{\mathcal D(P,U_2)}\le 1,
\qquad\text{i.e.}\qquad
\mathcal D(P,U_2)\le\frac{1}{9.49^2}\approx0.011.
\]
Hence for defects below approximately \(0.011\) the Lipschitz estimate
provides a non-trivial restriction on the min-entropy of a binary source.
\end{example}
\paragraph{Tsallis Entropy.}
The Tsallis entropy is
\[ T_\alpha(P) = \frac{1-\sum_i p_i^{\alpha}}{\alpha-1}.
\]
Its differentiability on the interior of the simplex again places it
within the smooth theory.
Hence it satisfies the same geometric stability estimates.
\paragraph{Collision Entropy.}
Collision entropy,
$H_2(P) = -\log \sum_i p_i^{2},$
is a special R\'enyi entropy of order two.
Accordingly, it inherits the same universal geometric estimates.
\paragraph{Jensen--Shannon Divergence.}
The Jensen--Shannon divergence is
\[ \operatorname{JSD}(P,Q) = \frac12
D_{\mathrm{KL}}(P\|M) + \frac12
D_{\mathrm{KL}}(Q\|M), \]
where
$M=\frac{P+Q}{2}.$
Since it is constructed from KL divergences, its smoothness immediately
implies the applicability of the universal Lipschitz and quadratic
bounds on compact subsets of the interior.
\subsection{General Smooth Functionals}\label{ch:sec:general-smooth-functionals}
The preceding examples illustrate a common principle.
\textbf{Principle.}  Every information functional satisfying
\(S\in C^{1}(\mathcal P)\) inherits the universal geometric stability
estimate
\[
\|S(P)-S(Q)\| \le L_S \sqrt{\mathcal D(P,Q)}.
\]
If \(S\in C^{2}(\mathcal P)\) and possesses a stationary point
\(\nabla_g S(P^{\star})=0\), then
\[
\|S(P)-S(P^{\star})\| \le C_S \mathcal D(P,P^{\star}).
\]
Thus the proofs become independent of the individual functional once
regularity has been verified.

\section{g-Entropy}\label{ch:sec:g-entropy}

The g-entropy \cite{briet2019properties} unifies Shannon, R\'enyi,
Tsallis \cite{tsallis1988possible}, min-entropy, and collision entropy under a single formula.
Let \(g:\mathbb R_{>0}\to\mathbb R\) be a continuous, strictly
monotone function and write \(h=g^{-1}\).  The g-entropy of a distribution \(P\) is

\[
H_g(P) = g\!\left(\sum_{i=1}^n P_i\, h(-\log P_i)\right).
\]

Different choices of \(g\) recover classical entropies:

\[
\begin{array}{ll}
g(x)=x & \text{Shannon entropy}, \\[2pt]
g(x)=\dfrac{\log x}{1-\alpha} & \text{R\'enyi entropy of order }\alpha, \\[2pt]
g(x)=\dfrac{x-1}{1-\alpha} & \text{Tsallis entropy}, \\[2pt]
g(x)=\lim_{\alpha\to\infty}\dfrac{\log x}{1-\alpha} & \text{Min-entropy}.
\end{array}
\]

Indeed, with \(g(x)=x\) one has \(h(y)=y\) and
\(\sum_i P_i(-\log P_i)=H_{\mathrm{Shannon}}(P)\); with
\(g(x)=\log(x)/(1-\alpha)\) one has \(h(y)=e^{(1-\alpha)y}\) and
\(\sum_i P_i e^{-(1-\alpha)\log P_i}=\sum_i P_i^\alpha\), so
\(H_g(P)=\log(\sum_i P_i^\alpha)/(1-\alpha)=H_\alpha(P)\).
The Tsallis case is identical with \(g(x)=(x-1)/(1-\alpha)\).

\begin{theorem}[g-Entropy Gradient]
\label{ch:thm:g-entropy-gradient}
Let \(h=g^{-1}\) and \(\Sigma_g(P)=\sum_i P_i h(-\log P_i)\).
The squared norm of the WIS gradient of the g-entropy is

\[
\|\nabla_g H_g(P)\|_g^2
=
g'(\Sigma_g(P))^2
\sum_{k=1}^n
P_k^2
\Bigl(
h(-\log P_k)-h'(-\log P_k)
-\sum_{j=1}^n P_j\bigl(h(-\log P_j)-h'(-\log P_j)\bigr)
\Bigr)^2.
\]

This expression is bounded on every compact subset of \(\Delta_n^\circ\).
\end{theorem}

\begin{proof}
The g-entropy is smooth on the interior simplex because \(g\) and
\(g^{-1}\) are smooth on their domains (all entropies of interest use
smooth \(g\)).  Writing \(\Sigma_g(P)=\sum_i P_i h(-\log P_i)\), the
chain rule gives
\[
\frac{\partial H_g}{\partial P_k}
=
g'(\Sigma_g)\Bigl(h(-\log P_k)-h'(-\log P_k)\Bigr),
\]
because \(\partial(P_k h(-\log P_k))/\partial P_k
= h(-\log P_k)-h'(-\log P_k)\) and the remaining terms
\(P_i h(-\log P_i)\), \(i\neq k\), are independent of \(P_k\).
In the logarithmic coordinates \(a_k=\log P_k\), the pullback metric
is the projected Euclidean metric, and the gradient of a function
\(S\) is represented by the components
\(b_k=P_k\,\partial S/\partial P_k-P_k\sum_j P_j\,\partial S/\partial P_j\)
(tangent to the simplex), whose squared norm is
\(\sum_k b_k^2-(1/n)(\sum_k b_k)^2\).  For \(S=H_g\) the second
term vanishes because \(\sum_k b_k\propto \sum_k P_k h(-\log P_k)
-\sum_j P_j h(-\log P_j)=0\), yielding the displayed formula.
Boundedness on compacts follows from continuity. \(\)
\end{proof}

\begin{corollary}[Universal g-Entropy Bound]
\label{ch:cor:g-entropy}
The g-entropy satisfies the Lipschitz estimate

\[
|H_g(P)-H_g(Q)|
\le
L_{g,K}\,\sqrt{\mathcal D(P,Q)}
\]

on any compact subset \(K\subset\Delta_n^\circ\), with Lipschitz constant
\(L_{g,K}\) equal to the supremum of the gradient norm over \(K\).
\end{corollary}

Thus the existence of a unified bound for all classical entropies is
not a coincidence: it follows from the g-entropy representation and the
universal geometry of the posterior manifold.
\section{Alpha-Divergences (Amari)}\label{ch:sec:alpha-divergences-amari}

The \(\alpha\)-divergences of Amari \cite{amari2016information} form a one-parameter family that
includes the Kullback--Leibler divergence as a limiting case and is
central to classical information geometry.  For
\(\alpha\in\mathbb R\setminus\{\pm1\}\),

\[
D_\alpha^{(A)}(P\|Q)
=
\frac{4}{1-\alpha^2}
\Bigl(
1-\sum_i P_i^{\frac{1-\alpha}{2}} Q_i^{\frac{1+\alpha}{2}}
\Bigr).
\]

For \(\alpha\to1\) this converges to the KL divergence
\(D_{\mathrm{KL}}(P\|Q)\), and for \(\alpha\to-1\) to
\(D_{\mathrm{KL}}(Q\|P)\).

\begin{theorem}[Alpha-Divergence is a Smooth Functional]
\label{ch:thm:alpha-div}
For every fixed \(\alpha\neq\pm1\) and every fixed reference
\(Q\in\Delta_n^\circ\), the map \(P\mapsto D_\alpha^{(A)}(P\|Q)\) is
smooth on the interior simplex.  Therefore the universal Lipschitz
estimate applies:

\[
|D_\alpha^{(A)}(P_1\|Q)-D_\alpha^{(A)}(P_2\|Q)|
\le
L_{\alpha,K}\,\sqrt{\mathcal D(P_1,P_2)},
\]

and if \(P_1=Q\) (the stationary point of the divergence), the
quadratic estimate

\[
D_\alpha^{(A)}(P\|Q)
\le
C_{\alpha,K}\,\mathcal D(P,Q)
\]

holds in a neighbourhood of \(Q\).
\end{theorem}

\begin{proof}
Smoothness follows from the fact that \(P\mapsto P_i^{\gamma}\) is
smooth on \((0,\infty)\) for any real exponent \(\gamma\).  The bounds
then follow from Theorems~\ref{ch:thm:14-1} and~\ref{ch:thm:14-2}.
The stationary point condition \(\nabla_g D_\alpha^{(A)}(\cdot\|Q)=0\)
at \(P=Q\) follows from the symmetry
\(D_\alpha^{(A)}(Q\|Q)=0\) and the first-order optimality condition.
\(\)
\end{proof}

The constant \(L_{\alpha,K}\) can be bounded explicitly by
differentiating the divergence in logarithmic coordinates.  Writing
\(a_i=\log(P_i/Q_i)\),

\[
\frac{\partial}{\partial a_i}D_\alpha^{(A)}(P\|Q)
=
\frac{2}{1+\alpha}
\,
\Bigl(
\frac{P_i}{Q_i}
\Bigr)^{\frac{1-\alpha}{2}}
-
\frac{2}{1-\alpha}
\,
\Bigl(
\frac{P_i}{Q_i}
\Bigr)^{-\frac{1+\alpha}{2}}.
\]

From this expression the gradient norm can be bounded uniformly on
compact subsets, yielding an explicit (if notationally heavy)
Lipschitz constant.

\subsection{Chi-Square Divergence}\label{ch:cor:chi-square}
The \(\chi^2\)-divergence is

\[
\chi^2(P\|Q) = \sum_i \frac{(P_i-Q_i)^2}{Q_i}.
\]

Although it diverges near the boundary of the simplex (when \(Q_i\to0\)
while \(P_i>0\)), it is smooth on every compact subset of the interior.

\begin{corollary}[Chi-Square Bound]

On any compact set \(K\subset\Delta_n^\circ\) where
\(\min_i Q_i\ge\varepsilon>0\),

\[
|\chi^2(P_1\|Q)-\chi^2(P_2\|Q)|
\le
L_{\chi^2,K}\,\sqrt{\mathcal D(P_1,P_2)},
\]

with \(L_{\chi^2,K}^2\le
(4/\varepsilon^2)\max_i P_i/(\min_i P_i)\).
\end{corollary}

\paragraph{Running example: browser fingerprinting.}
The $g$-entropy formulation applies directly to the browser fingerprinting example (Section~\ref{ch:sec:browser-fingerprinting-and-anonymity}).  The fingerprint distribution is strongly concentrated: a small number of browsers account for most users.  The $g$-entropy $H_g(P)=\sum_i g(P_i)$ with $g(p)=-p\log p$ (Shannon) captures this concentration, and Corollary~\ref{ch:cor:g-entropy} provides the universal quadratic bound.

\section{Hellinger Distance}\label{ch:sec:hellinger-distance}

The squared Hellinger distance

\[
H^2(P,Q) = 1-\sum_i\sqrt{P_iQ_i}
\]

is smooth on the entire interior simplex and bounded between 0 and 1.
Its WIS gradient is uniformly bounded, giving a global Lipschitz
constant independent of the compact set.

\begin{theorem}[Hellinger--UOD Comparison]
\label{ch:thm:hellinger}
For any \(P,Q\in\Delta_n^\circ\),

\[
H^2(P,Q) \le \frac14\,\mathcal D(P,Q).
\]

Moreover, for distributions in a neighbourhood of the uniform
distribution \(U_n\),

\[
\frac14\,\mathcal D(P,Q) - O\bigl(\mathcal D^{3/2}\bigr)
\le H^2(P,Q)
\le \frac14\,\mathcal D(P,Q).
\]

Thus the Hellinger distance is asymptotically equivalent to half the
UOD-geodesic distance near the uniform distribution.
\end{theorem}

\begin{proof}
Write \(P_i=e^{a_i}/n\) and \(Q_i=e^{b_i}/n\) in logarithmic coordinates.
Then

\[
\sqrt{P_iQ_i}=\frac1n e^{(a_i+b_i)/2}.
\]

Expanding the exponential to second order and summing gives the
quadratic approximation
\(H^2(P,Q)=\frac18\sum_i(a_i-b_i)^2+O(\|a-b\|^4)\).
The leading term equals \(\frac14\mathcal D(P,Q)\) by the definition
of the UOD (the Euclidean distance in the quotient space).  The
inequality follows from the concavity of the square root.
\(\)
\end{proof}

\section{Itakura--Saito Divergence}\label{ch:sec:itakura--saito-divergence}

The Itakura--Saito divergence \cite{sakamoto2008reversible} (used in audio processing and spectral
analysis) is a Bregman divergence generated by the convex potential
\(\phi(x)=-\log x\):

\[
D_{\mathrm{IS}}(P\|Q) = \sum_i
\Bigl(
\frac{P_i}{Q_i} - \log\frac{P_i}{Q_i} - 1
\Bigr).
\]

\begin{corollary}[Itakura--Saito Bound]
\label{ch:cor:is}
The Itakura--Saito divergence is smooth on every compact subset of
\(\Delta_n^\circ\).  By the Bregman theory of Chapter~\ref{ch:sec:universal-geometric-bounds}, it admits an
exact quadratic representation through its Average Hessian Operator,
and therefore satisfies

\[
\frac{\underline\lambda}{2}\,\mathcal D(P,Q)
\le
D_{\mathrm{IS}}(P\|Q)
\le
\frac{\overline\lambda}{2}\,\mathcal D(P,Q),
\]

where \(\underline\lambda\) and \(\overline\lambda\) are the extremal
eigenvalues of the Average Hessian of \(\phi(x)=-\log x\) along the
geodesic joining \(P\) and \(Q\).
\end{corollary}

\section{Beta-Divergences (Cichocki--Amari)}\label{ch:sec:beta-divergences-cichocki--amari}

The \(\beta\)-divergences \cite{cichocki2010families} form a
one-parameter family indexed by \(\beta\in\mathbb R\):

\[
D_\beta^{(B)}(P\|Q)
=
\frac{1}{\beta(\beta-1)}
\sum_i
\Bigl(
P_i^\beta
-
\beta P_i Q_i^{\beta-1}
+
(\beta-1)Q_i^\beta
\Bigr),
\]

where the cases \(\beta=0,1\) are defined by continuity and recover the
Itakura--Saito and Kullback--Leibler divergences, respectively.  For
\(\beta=2\) we obtain the squared Euclidean distance.

\begin{theorem}[Beta-Divergence Gradient]
\label{ch:thm:beta-div}
For \(\beta\neq0,1\) and a fixed reference \(Q\), the WIS gradient of
\(P\mapsto D_\beta^{(B)}(P\|Q)\) is

\[
(\nabla_g D_\beta^{(B)})_i
=
\frac{1}{\beta-1}
\,
\frac{
P_i^{\beta-1} - Q_i^{\beta-1}
}{
\sqrt{P_i}
}.
\]

The squared norm is bounded on every compact subset of
\(\Delta_n^\circ\).  Consequently, the universal Lipschitz estimate
applies with constant

\[
L_{\beta,K}^2
\le
\frac{1}{(\beta-1)^2}
\,
\frac{
\max_i\bigl(P_i^{\beta-1}-Q_i^{\beta-1}\bigr)^2
}{
\min_i P_i
}.
\]

At the stationary point \(P=Q\), the quadratic bound
\(D_\beta^{(B)}(P\|Q)\le C_{\beta,K}\,\mathcal D(P,Q)\) holds with

\[
C_{\beta,K}
\le
\frac{|\beta-1|}{2}
\,
\max_i Q_i^{\beta-2},
\]

obtained from the Hessian of the Bregman potential \(\phi(x)=x^\beta/\beta\).
\end{theorem}

\begin{proof}
Smoothness follows from the fact that \(x\mapsto x^\gamma\) is smooth on
\((0,\infty)\) for any real \(\gamma\).  The gradient is obtained by
differentiating the divergence and converting to logarithmic
coordinates.  Boundedness on compacts is immediate from continuity.
The quadratic constant follows from Theorem~\ref{ch:thm:14-2} and the
Hessian of \(\phi\). \(\)
\end{proof}

\section{R\'enyi Divergence}\label{ch:sec:renyi-divergence}

While Section~\ref{ch:sec:improved-bound-via-the-stationary-point} analysed the R\'enyi \emph{entropy}
\(H_\alpha(P)\), the R\'enyi divergence of order \(\alpha>0\),
\(\alpha\neq1\),

\[
D_\alpha(P\|Q)
=
\frac{1}{\alpha-1}
\log
\sum_i
P_i^\alpha Q_i^{1-\alpha},
\]

is a functional of two arguments.  When the second argument is fixed
(\(Q\) constant), it becomes a smooth functional of \(P\) alone.

\begin{theorem}[R\'enyi Divergence Bound]
\label{ch:thm:renyi-div}
For fixed \(\alpha>1\) and fixed reference \(Q\), the map
\(P\mapsto D_\alpha(P\|Q)\) is smooth on the interior simplex.  Its WIS
gradient is

\[
(\nabla_g D_\alpha)_i
=
\frac{\alpha}{\alpha-1}
\,
\frac{
P_i^{\alpha-1}Q_i^{1-\alpha}
-
\sum_j P_j^\alpha Q_j^{1-\alpha}
}{
\sqrt{\sum_j P_j^\alpha Q_j^{1-\alpha}}\;\sqrt{P_i}
}.
\]

The Lipschitz and quadratic bounds of
Theorems~\ref{ch:thm:14-1}--\ref{ch:thm:14-2} therefore apply with
constants depending on the compact set and on \(\alpha\).

At the reference point \(P=Q\), the divergence vanishes and its
second-order expansion in the logarithmic coordinates is
\[
D_\alpha(P\|Q)
=
\frac{\alpha}{2}\,\chi^2_Q(P,Q)
+
o\bigl(\chi^2_Q(P,Q)\bigr),
\qquad
\chi^2_Q(P,Q)=\sum_i \frac{(P_i-Q_i)^2}{Q_i}.
\]

Hence the leading term is proportional to the \emph{weighted}
chi-square divergence, not to the UOD: the two quadratic forms agree
only when the reference weights are balanced.  In particular, for the
uniform reference \(Q=U_n\), the identity
\(\chi^2_{U_n}(P,U_n)=\mathcal D(P,U_n)/n\) gives
\[
D_\alpha(P\|U_n)
=
\frac{\alpha}{2n}\,\mathcal D(P,U_n)
+
o\bigl(\mathcal D(P,U_n)\bigr).
\]

This gives an explicit quadratic constant
\(C_{\alpha,U_n}=\alpha/(2n)\) at the uniform stationary point,
consistent with the constant of
Remark~\ref{rem:improved-hessian}.
\end{theorem}

\begin{proof}
The gradient follows from differentiating the log-sum-exp expression.
For the second order at \(P=Q\), write \(P=Q+\delta\) with
\(\sum_i\delta_i=0\) and expand
\(\sum_i P_i^\alpha Q_i^{1-\alpha}\) to second order:
\(\sum_i Q_i(1+\delta_i/Q_i)^\alpha = 1 + \frac{\alpha(\alpha-1)}{2}\sum_i \delta_i^2/Q_i + O(\delta^3)\).
Taking the logarithm and multiplying by \(1/(\alpha-1)\) gives the
claimed leading term \((\alpha/2)\chi^2_Q(P,Q)\). \(\)
\end{proof}

\section{Max Leakage (Guessing Probability)}\label{ch:sec:max-leakage-guessing-probability}

The max leakage \cite{smith2009foundations, alvim2023information} of a
distribution relative to the uniform distribution is

\[
\mathcal L_\infty(P)
=
\log_2
\frac{
\max_i P_i
}{
1/n
}
=
\log_2 n - H_\infty(P).
\]

It is therefore proportional to the min-entropy shift and inherits its
piecewise-smooth structure.  The WIS gradient exists on every open
region where the maximizing index is unique and constant, and its norm
is given by

\[
\|\nabla_g\mathcal L_\infty(P)\|_g^2
=
\frac{
1
}{
\bigl(\min_i P_i\bigr)^2
}
\quad\text{(on the region where }\max_i P_i\text{ is attained at a
single index)}.
\]

\begin{theorem}[Max-Leakage Bound]
\label{ch:thm:max-leakage}
On any region where the maximizing index of \(P\) is unique and fixed,
the Lipschitz estimate

\[
|\mathcal L_\infty(P)-\mathcal L_\infty(Q)|
\le
\frac{1}{\min_i P_i}
\,
\sqrt{\mathcal D(P,Q)}
\]

holds, provided the maximizing index of \(Q\) coincides with that of
\(P\).  The quadratic estimate at the uniform distribution (which is a
stationary point of the max leakage) gives

\[
\mathcal L_\infty(P)
\le
\frac{1}{2(\min_i P_i)^2}
\,
\mathcal D(P,U_n).
\]

Across regions where the maximizing index changes, the global bound
remains the same, but the Lipschitz constant on the boundary may be
larger due to the non-smooth transition.
\end{theorem}

\begin{proof}
The gradient is obtained by differentiating
\(H_\infty(P)=-\log\max_i P_i\) on the regularity region.  The bounds
follow from Theorems~\ref{ch:thm:14-1} and~\ref{ch:thm:14-2} together
with the piecewise-smooth theory of Chapter~\ref{ch:sec:classical-smooth-information-functionals}. \(\)
\end{proof}

\section{Refined g-Entropy Bound}\label{ch:sec:refined-g-entropy-bound}

The gradient norm of the g-entropy computed in
Section~\ref{ch:sec:max-leakage-guessing-probability} can be made fully explicit for any differentiable
\(g\).  Writing \(h=g^{-1}\) and \(\Sigma_g(P)=\sum_i P_i h(-\log P_i)\),
Theorem~\ref{ch:thm:g-entropy-gradient} gives the exact squared norm

\[
\|\nabla_g H_g(P)\|_g^2
=
g'(\Sigma_g(P))^2
\sum_{k=1}^n
P_k^2
\Bigl(
A_k
-
\sum_{j=1}^n P_j A_j
\Bigr)^2,
\qquad
A_k=h(-\log P_k)-h'(-\log P_k).
\]

Since \(\sum_k P_k^2\le\max_i P_i\) on the simplex, bounding the
centred term by the maximum component gives the explicit upper bound

\[
\|\nabla_g H_g(P)\|_g^2
\le
g'(\Sigma_g(P))^2
\,
\max_k A_k^2
\,
\max_i P_i.
\]

\begin{corollary}[Lipschitz Constant for General g-Entropy]
\label{ch:cor:g-entropy-lipschitz}
For any \(g\) satisfying the smoothness hypotheses, the Lipschitz
constant on a compact set \(K\subset\Delta_n^\circ\) satisfies

\[
L_{g,K}^2
\le
\max_{P\in K}
g'(\Sigma_g(P))^2
\,
\max_{P\in K}\,\max_k A_k^2
\,
\max_{P\in K} P_i.
\]

For the special cases, applying the gradient formula of
Theorem~\ref{ch:thm:g-entropy-gradient} and bounding the projection term
by the maximum component gives:
\begin{align*}
g(x)=x &: \quad L_g^2 = \frac{1}{\min_i P_i},\\[4pt]
g(x)=\dfrac{\log x}{1-\alpha} &: \quad
L_\alpha^2 = \frac{\alpha^2}{(1-\alpha)^2}
\frac{\sum_i P_i^{2\alpha-2}}{(\sum_i P_i^\alpha)^2}
\max_i P_i,\\[4pt]
g(x)=\dfrac{x-1}{1-\alpha} &: \quad
L_{\mathrm{Tsallis}}^2 = \frac{\alpha^2}{(1-\alpha)^2}
\frac{\sum_i P_i^{2\alpha-2}}{(\sum_i P_i^\alpha)^2}
\max_i P_i.
\end{align*}
\end{corollary}

Thus the single g-entropy formula subsumes and unifies the Lipschitz
constants of all classical entropies, confirming that the independent
bounds derived in the previous sections are consequences of one common
structure.

\section{Summary of Explicit Constants}\label{ch:sec:summary-of-explicit-constants}

The following table summarises the Lipschitz and quadratic constants
for the information functionals treated in this chapter.

\[
\begin{array}{lll}
\textbf{Functional} & \textbf{Lipschitz } L_S^2 & \textbf{Quadratic } C_S \\ \hline
\text{Shannon entropy} & \displaystyle\frac{1}{\min_i P_i} & \displaystyle\frac{1}{2\min_i P_i} \\[8pt]
\text{R\'enyi entropy } \alpha & \displaystyle\frac{\alpha^2}{(1-\alpha)^2}\frac{\sum_i P_i^{2\alpha-2}}{(\sum_i P_i^{\alpha})^2}\max_i P_i & \displaystyle\frac{\alpha}{2|\alpha-1|}\frac{\max_i P_i^{\alpha-2}}{\sum_i P_i^{\alpha}} \\[8pt]
\text{Min-entropy} & \displaystyle\frac{\max_i P_i}{(\min_i P_i)^2} & \displaystyle\frac{1}{2(\min_i P_i)^2} \\[8pt]
\text{KL divergence} & \displaystyle\frac{1}{\min_i Q_i} & \displaystyle\frac{1}{2\min_i Q_i} \\[8pt]
\text{Chi-square} & \displaystyle\frac{4}{\varepsilon^2}\frac{\max_i P_i}{\min_i P_i} & \displaystyle\frac{2}{\varepsilon^2} \\[8pt]
\text{Hellinger} & \displaystyle\frac14 & \displaystyle\frac14
\end{array}
\]

All constants are evaluated on a compact set where
\(\min_i P_i\ge\varepsilon>0\) and \(\min_i Q_i\ge\varepsilon>0\) where
applicable.  The Hellinger constants are global.

\section{Empirical Validation and Implications}\label{ch:sec:empirical-validation-and-implications}

The bounds developed in this chapter are mathematical inequalities,
but their practical significance---and their limitations---are best
understood through concrete numerical experiments.  We present here
the results of a computational discovery campaign run on the
one-parameter family of one-heavy distributions

\[
P_\varepsilon = \bigl(1-(n-1)\varepsilon,\; \varepsilon,\; \ldots,\; \varepsilon\bigr),
\]

where $\varepsilon>0$ controls the concentration of probability mass on a single symbol.
This family smoothly interpolates between the uniform distribution
($\varepsilon=1/n$) and increasingly skewed configurations.

\paragraph{The Lipschitz bound is always tight but extremely loose.}
For every tested configuration, the Lipschitz estimate
$|H_\infty(P)-\log_2 n|\le L_\infty\sqrt{\mathcal D(P,U_n)}$ is
satisfied (as guaranteed by Theorem~\ref{ch:thm:14-1}).  However, the
margin is enormous.  Table~\ref{tab:empirical-lipschitz} shows
representative values.

\begin{table}[ht]
\centering
\caption{Empirical Lipschitz ratios for one-heavy distributions.}
\label{tab:empirical-lipschitz}
\begin{tabular}{c|c|c|c|c|c}
$n$ & $\varepsilon$ & $\mathcal D(P,U_n)$ & $L\sqrt{\mathcal D}$ & Actual gap & Ratio \\ \hline
$5$ & $0.005$ & $1.14$ & $42\,214$ & $2.29$ & $18\,400\times$ \\
$5$ & $0.100$ & $1.60$ & $98$ & $1.59$ & $62\times$ \\
$10$ & $0.005$ & $2.47$ & $61\,403$ & $3.26$ & $18\,800\times$ \\
$50$ & $0.005$ & $9.69$ & $108\,201$ & $5.24$ & $20\,700\times$ \\
$100$ & $0.005$ & $11.08$ & $94\,626$ & $5.66$ & $16\,700\times$
\end{tabular}
\end{table}

The Lipschitz ratio ranges from $62\times$ to over $20\,000\times$.
The bound is mathematically correct but quantitatively useless for most
practical purposes: the Lipschitz constant $L_\infty$ depends on
$1/(\min_i P_i)^2$, which explodes as any probability approaches zero.
This is not a defect of the proof but a fundamental limitation of the
linear (Lipschitz) approach: it captures only first-order information,
and near the boundary of the simplex the gradient norm diverges.

\paragraph{The quadratic bound is tighter but still loose for large $n$.}
The quadratic estimate $|H_2(P)-\log_2 n|\le C\,\mathcal D(P,U_n)$
with $C=1/n$ derived from the Hessian at the uniform distribution is
substantially tighter.  The ratio to the exact value
$H_2(P)=\log_2 n-D_2(P\|U_n)$ \cite{vanerven2014renyi} is

\[
\begin{array}{c|c|c|c|c|c}
n & \varepsilon & \mathcal D(P,U_n) & \text{WIS bound} & \text{Exact} & \text{Ratio} \\ \hline
5 & 0.05 & 2.56 & 0.51 & 1.18 & 2.3\times \\
5 & 0.005 & 1.14 & 0.23 & 1.57 & 6.9\times \\
10 & 0.05 & 2.96 & 0.30 & 1.18 & 3.9\times \\
10 & 0.005 & 2.47 & 0.25 & 2.21 & 8.8\times \\
50 & 0.01 & 8.04 & 0.16 & 2.58 & 16.1\times \\
50 & 0.005 & 9.69 & 0.19 & 3.35 & 17.3\times \\
50 & 0.001 & 4.56 & 0.09 & 3.81 & 41.8\times
\end{array}
\]

The WIS quadratic bound is $2\!-\!42\times$ wider than the exact
identity.  This ratio is smallest near the uniform distribution
(large $\varepsilon$, small $\mathcal D$) and grows with the dimension
$n$ and the skew.  The bound is therefore most informative precisely
where the exact R\'enyi divergence is also easiest to compute; for
highly skewed distributions in high dimensions the bound is a weak
estimate.

\paragraph{What these numbers imply.}
The empirical measurements confirm three structural facts about the
geometric bounds.

First, the Lipschitz bound is a qualitative, not quantitative, tool.
Its value lies in being universal---it applies to every smooth
functional without modification---not in being tight.  The Lipschitz
constant diverges near the simplex boundary; this is inevitable for any
linear estimate based on the gradient norm.  The quadratic bound is
quantitatively useful when the defect $\mathcal D(P,Q)$ is small, but
its constant can be improved by using the pointwise Hessian
$C_\alpha^\star(P)$ (Remark~\ref{rem:improved-hessian}) instead of the
uniform compact-set bound.

Second, the discovery campaign independently confirms the variational
structure established in Chapter~\ref{ch:sec:variational-uod}.  The extremal distributions found
by numerical optimisation---those that maximise the KL divergence from
the uniform distribution---collapse to at most three distinct
probability values, in accordance with the Three-Level Theorem
(Theorem~\ref{ch:thm:29-3}).  The numerical optimiser consistently returns
distributions with support size $3$--$7$ and a single dominant atom,
which is the two-level case of the theorem.

Third, the bounds derived in this chapter are not intended to replace
the specialised inequalities of information theory (Pinsker for KL \cite{csiszar2011information},
van Erven--Harremos for R\'enyi \cite{vanerven2014renyi}, etc.).  Their purpose is structural:
they show that all smooth information functionals, regardless of their
individual form, obey the same geometric stability principle.  The
constants may be loose, but the principle is universal.

\section{Discussion}\label{ch:sec:discussion-smooth_functionals}

The classical literature typically derives stability properties
separately for each entropy or divergence (see, e.g.,
\cite{csiszar2011information} for a classical treatment).
Within the WIS framework, these results are unified by the geometry of
the posterior manifold.
The Universal Optimality Defect provides the common quadratic reference
quantity, while the general theorems of Chapters~\ref{ch:sec:universal-geometric-bounds} and~\ref{ch:sec:the-average-hessian-operator-and-universal-bregman-theo} supply the
analytical machinery.
The next chapter extends this philosophy to non-smooth information
measures, including rank-based and maximum-based security functionals.
